\textbf{Phần 1. Trắc nghiệm lựa chọn}
\setcounterpageref{ex}{0}

\begin{ex} %[4D1V06]
Cho số $\overline{4a,7b2}$. Biết rằng số này chia hết cho 2 và khi làm tròn số đó đến hàng nghìn ta được số $46\,000$. Chữ số $a$ nhận giá trị là
\choice
{$4$}
{$6$}
{$7$}
{\True $5$}
\loigiai{
Xét số $\overline{4a,7b2}$ có chữ số hàng trăm là 7.\\
Vì $7 \ge 5$ nên khi làm tròn số đến hàng nghìn, ta phải làm tròn lên bằng cách cộng thêm 1 vào chữ số hàng nghìn $a$.\\
Theo bài ra ta được số $46\,000$, do đó:
$$a + 1 = 6 \implies a = 5$$
\textbf{Đáp án của câu hỏi là D}
}
\end{ex}

\begin{ex} %[4D1V06]
Cho biểu thức $P = a \times 25 + b \times 25 + c \times 25$. Biết $a + b + c = 400$. Giá trị của biểu thức $P$ là
\choice
{$8000$}
{\True $10\,000$}
{$12\,500$}
{$100\,000$}
\loigiai{
Áp dụng tính chất một số nhân với một tổng:
$$P = 25 \times (a + b + c)$$
Thay $a + b + c = 400$ vào biểu thức:
$$P = 25 \times 400 = 10\,000$$
\textbf{Đáp án của câu hỏi là B}
}
\end{ex}

\begin{ex} %[4D1V06]
Tổng của 5 số chẵn liên tiếp bằng $140$. Số lớn nhất trong 5 số chẵn đó là
\choice
{\True $32$}
{$30$}
{$34$}
{$28$}
\loigiai{
Vì có 5 số chẵn liên tiếp nên số ở chính giữa (số thứ ba) là:
$$140 : 5 = 28$$
Dãy 5 số chẵn liên tiếp đó là: $24; 26; 28; 30; 32$.\\
Số lớn nhất trong 5 số chẵn là $32$.\\
\textbf{Đáp án của câu hỏi là A}
}
\end{ex}

\begin{ex} %[4D1V06]
Tính giá trị của biểu thức sau bằng cách thuận tiện nhất:
$$48 \times 24 + 48 \times 75 + 48$$
\choice
{$4700$}
{$4752$}
{\True $4800$}
{$4848$}
\loigiai{
Ta viết số hạng 48 cuối cùng thành $48 \times 1$:
\begin{align*}
48 \times 24 + 48 \times 75 + 48 \times 1 &= 48 \times (24 + 75 + 1) \\
&= 48 \times 100 = 4800
\end{align*}
\textbf{Đáp án của câu hỏi là C}
}
\end{ex}

\begin{ex} %[4D1V06]
Một thửa đất hình chữ nhật có chu vi bằng $160\text{ m}$. Nếu giảm chiều dài đi $10\text{ m}$ và tăng chiều rộng thêm $10\text{ m}$ thì thửa đất trở thành một hình vuông. Diện tích thửa đất ban đầu là
\choice
{$1200\text{ m}^2$}
{$1400\text{ m}^2$}
{$1600\text{ m}^2$}
{\True $1500\text{ m}^2$}
\loigiai{
Nửa chu vi của thửa đất là:
$$160 : 2 = 80\text{ (m)}$$
Khi giảm chiều dài đi $10\text{ m}$ và tăng chiều rộng thêm $10\text{ m}$ thì hai cạnh bằng nhau, tức chiều dài hơn chiều rộng là:
$$10 + 10 = 20\text{ (m)}$$
Chiều dài ban đầu là:
$$(80 + 20) : 2 = 50\text{ (m)}$$
Chiều rộng ban đầu là:
$$50 - 20 = 30\text{ (m)}$$
Diện tích thửa đất ban đầu là:
$$50 \times 30 = 1500\text{ (m}^2\text{)}$$
\textbf{Đáp án của câu hỏi là D}
}
\end{ex}

\begin{ex} %[4D1V06]
Số lẻ lớn nhất có 5 chữ số khác nhau mà tổng các chữ số của nó bằng $25$ là
\choice
{\True $98\,521$}
{$98\,611$}
{$97\,621$}
{$98\,431$}
\loigiai{
Để số là lớn nhất, ta chọn các chữ số từ trái sang phải lớn nhất có thể:\\
- Hàng chục nghìn là 9, hàng nghìn là 8.\\
Tổng 3 chữ số còn lại là: $25 - (9 + 8) = 8$.\\
Vì số cần tìm là số lẻ nên hàng đơn vị phải lẻ: chọn hàng đơn vị là 1 để hàng trăm và chục lớn nhất.\\
Khi đó tổng chữ số hàng trăm và chục là: $8 - 1 = 7$.\\
Để số lớn nhất với các chữ số khác nhau, ta chọn hàng trăm là 5, hàng chục là 2.\\
Số tìm được là $98\,521$.\\
\textbf{Đáp án của câu hỏi là A}
}
\end{ex}

\begin{ex} %[4D1V06]
Một thư viện trong 3 tháng đầu năm mua sách mới. Tháng thứ nhất mua $1200$ cuốn sách. Tháng thứ hai mua số sách gấp 3 lần tháng thứ nhất. Tháng thứ ba mua số sách bằng trung bình cộng của cả ba tháng. Tháng thứ ba thư viện mua được số cuốn sách là
\choice
{$1800$ cuốn}
{$2000$ cuốn}
{\True $2400$ cuốn}
{$2800$ cuốn}
\loigiai{
Tháng thứ hai thư viện mua được số sách là:
$$1200 \times 3 = 3600\text{ (cuốn)}$$
Vì tháng thứ ba mua số sách bằng trung bình cộng của cả ba tháng nên số sách tháng thứ ba bằng trung bình cộng số sách của hai tháng đầu.\\
Tháng thứ ba thư viện mua được số sách là:
$$(1200 + 3600) : 2 = 2400\text{ (cuốn)}$$
\textbf{Đáp án của câu hỏi là C}
}
\end{ex}

\begin{ex} %[4D1V06]
Cho biểu thức $M = (x + y) : 2$. Với $x = 24\,500$ và $y = 35\,500$, giá trị của biểu thức $M$ là
\choice
{$25\,000$}
{\True $30\,000$}
{$35\,000$}
{$40\,000$}
\loigiai{
Thay các giá trị của $x$ và $y$ vào biểu thức:
\begin{align*}
M &= (24\,500 + 35\,500) : 2 \\
&= 60\,000 : 2 = 30\,000
\end{align*}
\textbf{Đáp án của câu hỏi là B}
}
\end{ex}

\begin{ex} %[4D1V06]
Một xưởng sản xuất nhập về $4$ cuộn thép, mỗi cuộn nặng $1250\text{ kg}$ và $6$ cuộn thép khác, mỗi cuộn nặng $1500\text{ kg}$. Tổng khối lượng thép xưởng đã nhập khi làm tròn đến hàng chục nghìn là
\choice
{$15\,000\text{ kg}$}
{$14\,000\text{ kg}$}
{$20\,000\text{ kg}$}
{\True $10\,000\text{ kg}$}
\loigiai{
Tổng khối lượng thép xưởng đã nhập là:
$$4 \times 1250 + 6 \times 1500 = 5000 + 9000 = 14\,000\text{ (kg)}$$
Vì chữ số hàng nghìn là 4 ($4 < 5$) nên khi làm tròn đến hàng chục nghìn ta làm tròn xuống thành $10\,000\text{ kg}$.\\
\textbf{Đáp án của câu hỏi là D}
}
\end{ex}

\begin{ex} %[4D1V06]
Có bao nhiêu số chẵn có 5 chữ số bắt đầu bằng chữ số $5$ và kết thúc bằng chữ số $6$?
\choice
{\True $1000$ số}
{$900$ số}
{$500$ số}
{$100$ số}
\loigiai{
Số có dạng $\overline{5abc6}$:\\
- Chữ số $a$ có 10 cách chọn (từ 0 đến 9).\\
- Chữ số $b$ có 10 cách chọn (từ 0 đến 9).\\
- Chữ số $c$ có 10 cách chọn (từ 0 đến 9).\\
Số lượng các số thỏa mãn là:
$$10 \times 10 \times 10 = 1000\text{ (số)}$$
\textbf{Đáp án của câu hỏi là A}
}
\end{ex}

\begin{ex} %[4D1V06]
Cho biểu thức $A = 20\,240 - x \times 5$. Tìm số tự nhiên $x$ để biểu thức $A$ đạt giá trị bé nhất lớn hơn 0.
\choice
{$x = 4048$}
{$x = 4046$}
{\True $x = 4047$}
{$x = 4040$}
\loigiai{
Để $A > 0$ thì $x \times 5 < 20\,240 \implies x < 20\,240 : 5 = 4048$.\\
Vì $x$ là số tự nhiên và để $A$ đạt giá trị bé nhất thì $x$ phải nhận giá trị lớn nhất thỏa mãn $x < 4048$, do đó $x = 4047$.\\
Khi đó $A = 20\,240 - 4047 \times 5 = 20\,240 - 20\,235 = 5$.\\
\textbf{Đáp án của câu hỏi là C}
}
\end{ex}

\begin{ex} %[4D1V06]
Để may $12$ bộ đồng phục học sinh cần $36\text{ m}$ vải. Hỏi nếu có $77\text{ m}$ vải cùng loại thì may được nhiều nhất bao nhiêu bộ đồng phục và còn thừa bao nhiêu mét vải?
\choice
{$24$ bộ, thừa $5\text{ m}$}
{\True $25$ bộ, thừa $2\text{ m}$}
{$25$ bộ, không thừa vải}
{$26$ bộ, thiếu $1\text{ m}$}
\loigiai{
Mỗi bộ đồng phục cần số mét vải là:
$$36 : 12 = 3\text{ (m)}$$
Thực hiện phép chia có dư:
$$77 : 3 = 25\text{ (dư } 2\text{)}$$
Vậy may được nhiều nhất 25 bộ đồng phục và còn thừa $2\text{ m}$ vải.\\
\textbf{Đáp án của câu hỏi là B}
}
\end{ex}

\textbf{Phần 2. Trắc nghiệm đúng sai}
\setcounterpageref{ex}{0}

\begin{ex} %[4D1V06]
Xét tính đúng sai của các khẳng định sau:
\choiceTF
{\True Số liền trước của số bé nhất có 5 chữ số là số lớn nhất có 4 chữ số.}
{\True Tích của ba số tự nhiên liên tiếp luôn là một số chẵn.}
{Giá trị của biểu thức $a \times b + c$ khi $a = 10, b = 20, c = 30$ bằng $500$.}
{\True Khi làm tròn số $49\,500$ đến hàng nghìn ta được số $50\,000$.}
\loigiai{
\begin{itemize}
    \item Số bé nhất có 5 chữ số là $10\,000$, số liền trước là $9999$ (số lớn nhất có 4 chữ số). Khẳng định a Đúng.
    \item Trong ba số tự nhiên liên tiếp luôn có ít nhất một số chẵn, do đó tích của ba số này luôn là số chẵn. Khẳng định b Đúng.
    \item Thay số: $10 \times 20 + 30 = 200 + 30 = 230 \neq 500$. Khẳng định c Sai.
    \item Chữ số hàng trăm là 5 nên khi làm tròn đến hàng nghìn ta làm tròn lên thành $50\,000$. Khẳng định d Đúng.
\end{itemize}
\begin{itemize}
    \item a) \textbf{Đúng}.
    \item b) \textbf{Đúng}.
    \item c) \textbf{Sai}.
    \item d) \textbf{Đúng}.
\end{itemize}
}
\end{ex}

\begin{ex} %[4D1V06]
Một trang trại nuôi $2400$ con gà. Số con vịt bằng một nửa số con gà. Số con ngan ít hơn số con vịt là $300$ con. Xét tính đúng sai của các khẳng định sau:
\choiceTF
{\True Trang trại có $1200$ con vịt.}
{\True Trang trại có $900$ con ngan.}
{Tổng số con gà và con vịt của trang trại là $3800$ con.}
{\True Tổng số gia cầm của cả trang trại là $4500$ con.}
\loigiai{
\begin{itemize}
    \item Số vịt: $2400 : 2 = 1200\text{ (con)}$. Khẳng định a Đúng.
    \item Số ngan: $1200 - 300 = 900\text{ (con)}$. Khẳng định b Đúng.
    \item Tổng gà và vịt: $2400 + 1200 = 3600\text{ (con)}$ (không phải $3800$ con). Khẳng định c Sai.
    \item Tổng số gia cầm: $2400 + 1200 + 900 = 4500\text{ (con)}$. Khẳng định d Đúng.
\end{itemize}
\begin{itemize}
    \item a) \textbf{Đúng}.
    \item b) \textbf{Đúng}.
    \item c) \textbf{Sai}.
    \item d) \textbf{Đúng}.
\end{itemize}
}
\end{ex}

\textbf{Phần 3. Trả lời ngắn}
\setcounterpageref{ex}{0}

\begin{ex} %[4D1V06]
Tính tổng của số chẵn lớn nhất có 5 chữ số khác nhau và số lẻ bé nhất có 5 chữ số khác nhau.
\shortans{108999}
\loigiai{
\textbf{Đáp số:} $108999$\\
- Số chẵn lớn nhất có 5 chữ số khác nhau là $98\,764$.\\
- Số lẻ bé nhất có 5 chữ số khác nhau là $10\,235$.\\
Tổng của hai số đó là:
$$98\,764 + 10\,235 = 108\,999$$
}
\end{ex}

\begin{ex} %[4D1V06]
Tìm giá trị của $x$, biết: $(x + 2500) \times 4 = 50\,000$.
\shortans{10000}
\loigiai{
\textbf{Đáp số:} $10000$\\
Ta có:
\begin{align*}
(x + 2500) \times 4 &= 50\,000 \\
x + 2500 &= 50\,000 : 4 \\
x + 2500 &= 12\,500 \\
x &= 12\,500 - 2500 \\
x &= 10\,000
\end{align*}
}
\end{ex}

\begin{ex} %[4D1V06]
Hiệu của hai số bằng $36$. Nếu tăng số bị trừ thêm $14$ đơn vị và giảm số trừ đi $6$ đơn vị thì hiệu mới của hai số bằng bao nhiêu?
\shortans{56}
\loigiai{
\textbf{Đáp số:} $56$\\
Khi tăng số bị trừ thêm 14 đơn vị thì hiệu tăng 14 đơn vị.\\
Khi giảm số trừ đi 6 đơn vị thì hiệu tăng tiếp 6 đơn vị.\\
Hiệu mới của hai số là:
$$36 + 14 + 6 = 56$$
}
\end{ex}

\begin{ex} %[4D1V06]
Một hình vuông có chu vi bằng chu vi hình chữ nhật có chiều dài $28\text{ cm}$ và chiều rộng $16\text{ cm}$. Tính diện tích của hình vuông đó theo đơn vị xăng-ti-mét vuông.
\shortans{484}
\loigiai{
\textbf{Đáp số:} $484$\\
Chu vi hình chữ nhật là:
$$(28 + 16) \times 2 = 88\text{ (cm)}$$
Cạnh của hình vuông là:
$$88 : 4 = 22\text{ (cm)}$$
Diện tích hình vuông là:
$$22 \times 22 = 484\text{ (cm}^2\text{)}$$
}
\end{ex}

\begin{ex} %[4D1V06]
Tìm số tự nhiên có 3 chữ số, biết rằng khi viết thêm chữ số $2$ vào bên trái số đó thì ta được số mới gấp $9$ lần số ban đầu.
\shortans{250}
\loigiai{
\textbf{Đáp số:} $250$\\
Gọi số cần tìm có dạng $\overline{abc}$ ($100 \le \overline{abc} \le 999$).\\
Khi viết thêm chữ số 2 vào bên trái ta được số $2000 + \overline{abc}$.\\
Theo bài ra ta có:
\begin{align*}
2000 + \overline{abc} &= 9 \times \overline{abc} \\
8 \times \overline{abc} &= 2000 \\
\overline{abc} &= 2000 : 8 \\
\overline{abc} &= 250
\end{align*}
Vậy số cần tìm là $250$.
}
\end{ex}

\begin{ex} %[4D1V06]
Tính giá trị của biểu thức bằng cách thuận tiện nhất:
$$125 \times 7 \times 8$$
\shortans{7000}
\loigiai{
\textbf{Đáp số:} $7000$\\
Áp dụng tính chất giao hoán và kết hợp:
\begin{align*}
125 \times 7 \times 8 &= (125 \times 8) \times 7 \\
&= 1000 \times 7 = 7000
\end{align*}
}
\end{ex}
